% id: 115-17
% book: 베이직쎈 확률과 통계 · 06 이항분포와 정규분포
% section: 유형 071 정규분포의 활용; 확률 구하기
% figure: tikz:fig-115-17
% answer: 
% answer_source: 
% variant_level: 0 (원본 전사)
% 이 파일은 build.mjs 가 items.json 에서 만든다. 고칠 때는 items.json 을 고친다.
\smash{\raisebox{1.5mm}{\dmkichul{베이직쎈 확률과 통계 115쪽 17번}}}
\noindent\begin{minipage}[t]{0.58\linewidth}\vspace{0pt}\raggedright 어느 가게에서 판매하는 사탕 한 봉지에 들어 있는 사탕의 개수는 평균이 58, 표준편차가 6인 정규분포를 따른다고 한다. 오른쪽 표준정규분포표를 이용하여 사탕의 개수가 52 이상 67 이하인 봉지는 전체의 몇 $\%$인지 구하면?\end{minipage}\hfill\begin{minipage}[t]{0.4\linewidth}\vspace{0pt}\centering \resizebox{\ifdim\width>\linewidth\linewidth\else\width\fi}{!}{% 115-17(인쇄 115쪽) — 115-17 오른쪽에 놓인 표준정규분포표($z$ / $\mathrm{P}(0\le Z\le z)$).
% 스캔 표라 크롭하지 않고 tabular 로 다시 조판했다(칸 값은 원문 그대로 · 원문에 없는 행은 넣지 않는다).
{\setlength{\tabcolsep}{5pt}\renewcommand{\arraystretch}{1.25}%
\begin{tabular}{|c|c|}
\hline
$z$ & $\mathrm{P}(0\le Z\le z)$ \\
\hline
$0.5$ & $0.19$ \\
\hline
$1.0$ & $0.34$ \\
\hline
$1.5$ & $0.43$ \\
\hline
\end{tabular}}
}\end{minipage}\par
\dmchoicesauto{}{$38\mkern3mu \%$}{$53\mkern3mu \%$}{$68\mkern3mu \%$}{$77\mkern3mu \%$}{$86\mkern3mu \%$}
